\scanpage{16}
\uline{opérant} \deleted{sans points fixes} \inserted{\uline{librement}} sur $X_{pq}$, en \unsure{prenant le|passant au} quotient de $C_{pq}$ par $\pi$ (qui s'identifie au groupe fondamental de $X\ldots$) Elles sont donc munies (canoniquement modulo isotopie de la carte) d'une métrique riemannienne (la métrique quotient) pour laquelle la carte est \uline{géodésique}. \deleted{\unsure{et}} Elle est régulière sss le sous-groupe $\pi$ de $\Gamma_{pq}$ est invariant.
% slanted
\margin{\uuline{attention} au cas où $p$ ou $q$ est infini ! \unsure{(passant|passant} au quotient, il peut y avoir création de nouveaux sommets comb. d'ordre fini\ldots}
% slanted
\margin{Notion de carte \uline{métriquement} quasi-régulière (\unsure{complète|complétée} \deleted{\unsure{et quasi-rég}}, i.e. dont le rev. universel est carte géod. standart) \unsure{Donner|Donnes} encore th. d'équiv. de cat. (\uline{sans} \unsure{repèr|répèr} \ldots)}

\uline{NB} Il y a donc lieu, pour \inserted{tt} $(p,q)$ donné, pour déterminer les $\pi$ possibles, d'étudier \deleted{\unsure{quels}} l'op. de $\Gamma_{pq}$ sur $X_{pq}$ -- et plus particulièrement, de trouver les \deleted{pts de $X_{pq}$} \deleted{\unsure{groupes de}} stabilisateurs dans $\Gamma_{pq}$ des pts de $X_{pq}$.

\deleted{\unsure{Corollaire}} \uline{Re NB} \unsure{Si|Il} \deleted{\unsure{on a|y a}} la carte \inserted{qu. rég.} est orientée \inserted{(p.ex. \uline{\unsure{orient.bl|orient.ble|orientable}} et épinglée)}, \deleted{\unsure{elle}} son représentant métriquement quasi-régulier se trouve muni canoniquement d'une structure \uline{complexe} -- qui, dans le cas compact \inserted{(i.e. carte \uline{finie})} \inserted{(i.e. $p,q$ finis et $\pi$ d'\uline{indice} \uline{fini})} en fait donc une \uline{courbe algébrique} projective \unsure{de ma suite|de ma récolte|de marsault|de ma seule} ! \unsure{Sa} métrique est celle de Bergman.
% archive sheet number: 15

